\documentclass[12pt,a4paper]{article}
\usepackage[top=22mm,bottom=24mm,left=27mm,right=27mm]{geometry}
\usepackage{amsmath,amssymb,xcolor}
\definecolor{link}{HTML}{205973}
\usepackage[hypertexnames=false,colorlinks=true,linkcolor=link,urlcolor=link]{hyperref}
\hypersetup{pdftitle={Theorem 2. A weighted approximation guarantee for the sampled value function}}
\setlength{\parindent}{0pt}
\setlength{\parskip}{5pt}
\raggedbottom
\newcommand{\R}{\mathbb R}
\newcommand{\E}{\mathbb E}
\newcommand{\Prob}{\mathbb P}
\newcommand{\lead}[1]{\par\smallskip\noindent\textbf{#1}\enspace}
\newcommand{\piopt}{\pi_{\mu^*,\nu}}
\newcommand{\class}{\mathcal J}
\newcommand{\wn}[1]{\left\|#1\right\|_{\infty,1/\psi}}
\begin{document}
{\LARGE\bfseries Theorem 2. A weighted approximation\\[2pt] guarantee for the sampled value function\par}
\medskip
A large approximation error at a rarely visited state can dominate the usual
sup-norm bound. Theorem~2 compares the same fitted value function against
every state weighting. How much a weight helps depends on both the states
visited and how much the weight can grow in one transition.

\section*{The model and the guarantee}
Let \(\mathcal S,\mathcal A\) be finite nonempty state and action sets.
In state \(x\), action \(a\) costs \(g_{x,a}\in\R\) and moves to \(y\) with
probability \(p(x,y,a)\). For discount \(0<\alpha<1\), let
\[
 (TJ)(x)=\min_{a\in\mathcal A}
       \Bigl\{g_{x,a}+\alpha\sum_y p(x,y,a)J(y)\Bigr\}.
\]
The optimal discounted cost \(J^*\) satisfies \(TJ^*=J^*\).
Fix an optimal policy \(\mu^*\) and an initial distribution \(\nu\).
Draw \(N\ge1\) states \(X_1,\ldots,X_N\) independently from the discounted
occupancy \(\piopt\) defined on page~\pageref{eq:occupancy}.
Write \(\rho^\top f=\sum_x\rho_x f(x)\) and
\(\|f\|_{1,\nu}=\sum_x\nu_x|f(x)|\).

Let \(\Phi:\mathcal S\to\mathcal H\) take values in a real Hilbert space.
The class \(\class\) consists of \(\widetilde J_{\mathbf z,b}(x)
=\langle\Phi(x),\mathbf z\rangle+b\) with \(\|\mathbf z\|\le C\), \(|b|\le B\),
and \(C,B\ge0\). Fit \(\widehat J\) by solving
\begin{equation}
 \begin{aligned}
  \operatorname{maximize}_{J\in\class,\,s\in\R^N}\quad&
      \nu^\top J-\frac{2}{(1-\alpha)N}\sum_{i=1}^N s_i\\
  \text{subject to}\quad& J(X_i)\le(TJ)(X_i)+s_i,\quad s_i\ge0,\quad 1\le i\le N .
 \end{aligned}
 \tag{18}
\end{equation}
For each weight \(\psi:\mathcal S\to[1,\infty)\), define
\[
 \wn{f}=\max_x\frac{|f(x)|}{\psi(x)},\qquad
 \beta(\psi)=\max_{x,a}\frac{\sum_y p(x,y,a)\psi(y)}{\psi(x)} .
\]
Write \(K=\max_x\|\Phi(x)\|\), \(\|g\|_\infty=\max_{x,a}|g_{x,a}|\),
and set \(D=4CK(1+\alpha)+4B(1-\alpha)+2\|g\|_\infty\).

\textbf{Theorem 2.} Let \(\epsilon>0\), \(0<\delta<1\) and
\(N\ge D^2(1+\sqrt{2\ln(1/\delta)})^2/\epsilon^2\).
With probability at least \(1-\delta-\delta^4\), every optimizer of (18) satisfies
\begin{equation}
 \|J^*-\widehat J\|_{1,\nu}
 \le \inf_{\substack{V\in\class\\ \psi\ge1}}\!
 \wn{J^*-V}
 \left(\nu^\top\psi+
 \frac{2(\piopt^\top\psi)(1+\alpha\beta(\psi))}{1-\alpha}\right)
 +\frac{4\epsilon}{1-\alpha}.
 \tag{25}\label{eq:theorem}
\end{equation}

\newpage
\section*{What the weight controls}
The weighted norm scales the error at \(x\) by \(1/\psi(x)\). The term
\(\nu^\top\psi\) accounts for this scaling at the starting state.
The term \(\piopt^\top\psi\) accounts for it over discounted visits,
and \(\beta(\psi)\) measures the largest expected one-step growth in the
weight. A weight is used only to evaluate the bound. The optimization
in (18) is the same for every choice.

For a policy \(\mu\), let \(P_\mu(x,y)=p(x,y,\mu(x))\),
\(g_\mu(x)=g_{x,\mu(x)}\), and \(T_\mu J=g_\mu+\alpha P_\mu J\).
Its discounted state distribution is
\begin{equation}
 \pi_{\mu,\nu}^\top=(1-\alpha)\sum_{t=0}^\infty\alpha^t\nu^\top P_\mu^t
       =(1-\alpha)\nu^\top(I-\alpha P_\mu)^{-1}.
 \tag{17}\label{eq:occupancy}
\end{equation}
Each \(\nu^\top P_\mu^t\) is a probability row and the mixing weights sum
to one. The matrix series converges entrywise because all entries of
\(P_\mu^t\) lie in \([0,1]\). Multiplying its partial sums by
\(I-\alpha P_\mu\) gives the inverse formula by telescoping.

\lead{A weighted Bellman estimate.}
Write \(d=\wn{J^*-V}\) and \(r(V)=(V-TV)_+\), where \(u_+=\max\{u,0\}\)
coordinatewise. For every \(x,a\), the one-step action values for \(V\)
and \(J^*\) differ by at most
\[
 \alpha\sum_y p(x,y,a)|V(y)-J^*(y)|
 \le \alpha d\sum_y p(x,y,a)\psi(y)
 \le \alpha\beta(\psi)d\,\psi(x).
\]
Taking the minimum over actions preserves the bound on the difference.
Using \(TJ^*=J^*\) gives \(r(V)(x)\le(1+\alpha\beta(\psi))d\,\psi(x)\).
This is a one-step estimate and holds for every \(\psi\ge1\), even when
\(\alpha\beta(\psi)\ge1\).

\lead{Choose one comparator before drawing the sample.}
Let \(P f=\piopt^\top f\) and define the population criterion
\[
 L(V)=\nu^\top(J^*-V)+\frac{2}{1-\alpha}\,P r(V).
\]
Averaging \(J^*-V\le d\psi\) under \(\nu\) and the residual bound under
\(\piopt\) shows that \(L(V)\) is at most the product under the infimum
in \eqref{eq:theorem}, for each \(V,\psi\).
There is a deterministic minimizer \(V_0\) of \(L\) on \(\class\).
Indeed, projecting \(\mathbf z\) onto the finite span of the
\(\Phi(x)\) preserves every fitted value and weakly decreases its norm.
Thus \(\class\) is the continuous image of a compact ball times \([-B,B]\).
The Bellman operator is a finite minimum of affine maps, so \(L\) is
continuous. It follows that \(L(V_0)\) is bounded by the joint infimum
in \eqref{eq:theorem}. This argument allows the infimum over weights to
be unattained.

\newpage
\section*{From a sampled optimum to a population bound}
Write \(P_N f=N^{-1}\sum_i f(X_i)\) and \(r_\mu(J)=(J-T_\mu J)_+\).
For a fixed \(J\), the least feasible slack at draw \(i\) in (18) is
\(r(J)(X_i)\). Substitution gives the continuous reduced objective
\(\nu^\top J-2P_Nr(J)/(1-\alpha)\). Compactness of \(\class\) gives
a maximizer on every draw, and every optimizer \(\widehat J\) satisfies
\begin{equation}
 \nu^\top(J^*-\widehat J)+\frac{2}{1-\alpha}P_Nr(\widehat J)
 \le \nu^\top(J^*-V_0)+\frac{2}{1-\alpha}P_Nr(V_0).
 \tag{a}\label{eq:opt}
\end{equation}

For the optimal policy, \(T_{\mu^*}J^*=J^*\). Subtracting its Bellman
equation from \(J-T_{\mu^*}J\) gives
\[
 J-J^*=(I-\alpha P_{\mu^*})^{-1}(J-T_{\mu^*}J).
\]
The inverse has nonnegative entries by (17). Replacing the residual on
the right by its positive part then bounds \((J-J^*)_+\). Averaging and
using \(|u|=-u+2u_+\) at \(u=J(x)-J^*(x)\) gives
\begin{equation}
 \|J^*-J\|_{1,\nu}
 \le\nu^\top(J^*-J)+\frac{2}{1-\alpha}P r_{\mu^*}(J).
 \tag{b}\label{eq:loss}
\end{equation}

\lead{Two sampling events.}
We need the following two inequalities on an event of probability at
least \(1-\delta-\delta^4\).
\begin{align}
 P r_{\mu^*}(J)&\le P_N r_{\mu^*}(J)+\epsilon
           &&\text{for every }J\in\class, \tag{c}\label{eq:uniform}\\
 P_Nr(V_0)&\le P r(V_0)+\epsilon
           &&\text{for the fixed }V_0. \tag{d}\label{eq:fixed}
\end{align}
The directions differ because the fitted function and the comparator are on
opposite sides of \eqref{eq:opt}. In \eqref{eq:loss}, take \(J=\widehat J\).
Apply \eqref{eq:uniform} and use \(r_{\mu^*}(\widehat J)\le r(\widehat J)\),
which follows from \(T\widehat J\le T_{\mu^*}\widehat J\).
Then \eqref{eq:opt} and \eqref{eq:fixed} give
\begin{align*}
 \|J^*-\widehat J\|_{1,\nu}
 &\le \nu^\top(J^*-V_0)+\frac{2}{1-\alpha}\bigl(P_Nr(V_0)+\epsilon\bigr)\\
 &\le L(V_0)+\frac{4\epsilon}{1-\alpha}.
\end{align*}
The two deviations each contribute \(2\epsilon/(1-\alpha)\).
Since \(L(V_0)\) is at most the infimum in \eqref{eq:theorem}, these
inequalities prove the theorem on the same event for every optimizer.

\lead{The concentration inequality used below.}
For independent inputs, if changing coordinate \(i\) changes a statistic
\(Z\) by at most \(c_i\), then for \(t>0\) and \(\sum_i c_i^2>0\),
\[
 \Prob\{Z-\E Z\ge t\}\le\exp\!\left(-\frac{2t^2}{\sum_i c_i^2}\right).
\]
Indeed, exposing the inputs one by one gives conditional-mean increments
with mean zero and conditional range of width at most \(c_i\). For any
centered variable in an interval of width \(c\), the second derivative
of its log moment-generating function is a tilted variance, at most
\(c^2/4\). Integrating twice bounds that log by \(\lambda^2 c^2/8\).
Iterated conditioning and Markov's inequality at
\(\lambda=4t/\sum_i c_i^2\) give the stated bound.

\newpage
\section*{The two probability estimates}
For each fixed policy write \(d_\mu(x)=\Phi(x)-\alpha\sum_yP_\mu(x,y)\Phi(y)\).
Its norm is at most \(K(1+\alpha)\) by the triangle inequality, and
\[
 r_\mu(\widetilde J_{\mathbf z,b})(x)
   =\bigl(\langle d_\mu(x),\mathbf z\rangle+(1-\alpha)b-g_\mu(x)\bigr)_+ .
\]
Consequently \(0\le r_\mu(J)\le CK(1+\alpha)+B(1-\alpha)+\|g\|_\infty\le D/2\).
Taking the maximum over actions gives the same range for \(r(J)\).

\lead{Uniform deviation for the fit.}
Fix the sample and let \(\sigma_i\) be independent uniform signs in
\(\{-1,1\}\). For the class \(\mathcal F=\{r_{\mu^*}(J):J\in\class\}\), set
\[
 \widehat R_N(\mathcal F)=
 \E_\sigma\sup_{f\in\mathcal F}\frac1N\sum_i\sigma_i f(X_i).
\]
A coordinatewise 1-Lipschitz transformation cannot increase this quantity.
Indeed, condition on all signs except one. For each candidate \(j\), let
\(a_j\) be its sum over the other coordinates and \(u_j\) its remaining
score, before division by \(N\). Twice the conditional expected supremum of the
unnormalized sum is
\(\sup_{j,k}(a_j+a_k+\varphi(u_j)-\varphi(u_k))\).
Replacing the difference by \(|u_j-u_k|\) bounds it above. Exchanging
\(j,k\) identifies the bound with the untransformed expression. Iterate
over coordinates with \(\varphi_x(u)=(u-g_{\mu^*}(x))_+\).

For the linear term write \(d_i=d_{\mu^*}(X_i)\).
Its signed supremum is at most \(C\|\sum_i\sigma_i d_i\|/N\).
The signs give \(\E_\sigma\|\sum_i\sigma_i d_i\|^2=\sum_i\|d_i\|^2\)
and \(\E_\sigma|\sum_i\sigma_i|\le\sqrt N\). Applying Cauchy--Schwarz
to the linear term and adding the intercept contribution gives
\[
 2\widehat R_N(\mathcal F)
 \le \frac{2CK(1+\alpha)+2B(1-\alpha)}{\sqrt N}\le\frac{D}{\sqrt N}.
\]
For \(Z=\sup_{f\in\mathcal F}(Pf-P_N f)\), an independent sample with mean
\(P_N'\) gives \(\E Z\le\E\sup_f(P_N'f-P_N f)\). Swapping each pair of
observations independently inserts a random sign. Splitting the difference
into two signed sums yields \(\E Z\le2\E\widehat R_N(\mathcal F)\).
All functions of the finite sample vector are measurable.

Suppose \(D>0\). Changing one observation changes \(Z\) by at most
\(2D/N\) since \(|f|\le D\). The bounded-differences inequality gives
\[
 Z\le\frac{D}{\sqrt N}+D\sqrt{\frac{2\ln(1/\delta)}{N}}\le\epsilon
 \qquad\text{with probability at least }1-\delta.
\]
This proves \eqref{eq:uniform}.

\lead{Fixed-comparator deviation.}
For \(D>0\), the independent variables \(r(V_0)(X_i)\) lie in \([0,D/2]\).
Changing one input changes their mean by at most \(D/(2N)\). The same
inequality now bounds the failure probability in \eqref{eq:fixed} by
\(\exp(-8N\epsilon^2/D^2)\le\delta^4\), since
\(N\epsilon^2/D^2\ge2\ln(1/\delta)\). If \(D=0\) every residual is zero.
A union bound proves the claimed probability. \hfill\(\square\)

\newpage
\section*{A two-state example with no best weight}
There are two absorbing states \(s,t\), one action, discount \(\alpha=1/2\),
and stage costs \(g(s)=0\), \(g(t)=5\). Begin with distribution
\(\nu=(9/10,1/10)\). Since the state never changes, \(\piopt=\nu\).
The costs to go are \(J^*(s)=0\) and \(J^*(t)=5/(1-1/2)=10\).
Take \(\Phi=0\), \(C=B=0\), so the only admissible fit is \(V=0\).
Its actual \(\nu\)-weighted error is \(1\).

Let the weight have values \(\psi(s)=u,\ \psi(t)=v\) with \(u,v\ge1\).
Every transition preserves its state's weight, giving \(\beta(\psi)=1\).
Hence the approximation term in (25), before the additive sampling error, is
\[
 \frac{10}{v}\left(\frac{9u+v}{10}\right)
       \left(1+\frac{2(1+1/2)}{1-1/2}\right)
       =7+63\,\frac{u}{v}.
\]
\begin{center}
\begin{tabular}{lcc}
Weight \(\bigl(\psi(s),\psi(t)\bigr)\) &
\(\wn{J^*-V}\) & Approximation term\\[3pt]\hline
\((1,1)\) & \(10\) & \(70\)\\
\((1,10)\) & \(1\) & \(13.3\)\\
\((1,v)\) as \(v\to\infty\) & \(10/v\to0\) & \(7+63/v\to7\)
\end{tabular}
\end{center}
Every finite admissible weight gives a value strictly larger than \(7\).
Their infimum is exactly \(7\). The same value function is fitted in
all three rows.

\section*{Relation to the original result}
Equation (25) holds with the published constants and the joint infimum.
Appendix B.2 reverses the sign of the comparator slack in one step.
Feasibility of (18) requires \((V-TV)_+\). The fixed-comparator
probability estimate also depends on which comparator is chosen.
Minimizing \(L\) before sampling justifies the infimum on one event,
including the nonattainment illustrated above.

For \(\psi\equiv1\), both averages of the weight and its drift equal one.
The coefficient in (25) reduces to \((3+\alpha)/(1-\alpha)\), recovering
the printed Theorem~1. The direct unweighted proof bounds the comparator
residual at every sampled state and therefore saves its sampling error.
The stated range is \(0<\delta<1\). For \(\delta\ge1\),
the printed probability lower bound is nonpositive.

{\small\textbf{Sources.}
\href{https://moallemi.com/ciamac/papers/kernel-adp-2012.pdf\#page=28}{Bhat, Farias, Moallemi and Zheng, Theorem~2 and Appendix~B}.
\href{https://github.com/nikhilpb/kadp/blob/main/Kadp/Theorem2.lean}{Formal theorem},
\href{https://github.com/nikhilpb/kadp/blob/main/Kadp/Weighted.lean}{weighted Bellman estimate},
\href{https://github.com/nikhilpb/kadp/blob/main/Kadp/PopulationComparison.lean}{sampling comparison},
and \href{https://github.com/nikhilpb/kadp/blob/main/Kadp/Examples/Theorem2.lean}{two-state example}.}
\end{document}
